\section{Introduction}
\label{sec:introduction}
\subsection{Motivation}
Algebraic tensor rewriting is used to optimize tensor computations. A rewrite
rule can be simple in algebraic form, but complex to implement. For example,
dot reassociation changes $(XY)Z$ into $X(YZ)$. The equation shows how to
rearrange the computation, but general tensor dot operations also need
attributes specifying their batch and contraction dimensions. The rewrite
must compute these attributes for the new operations while preserving the
output axis order. XLA's \code{AssociativeReorderNestedDot} routine illustrates
this work: it includes dimension-mapping logic as well as checks on when the
rewrite applies and whether it is beneficial~\cite{xla-simplifier}.

Computing these replacement attributes is a major source of complexity.
Matching the operations and attributes on the left-hand side identifies the
source computation, but does not determine the attributes needed on the
right-hand side. Users must write additional expressions or custom code to
compute them. This requires tracking how axis positions, operand pairings,
and output order change across composed operations. Two axes with the same
size can still represent different coordinates, so sizes alone are not enough
to determine their correspondence. The transformation must also work for
tensors of arbitrary rank; checking a few fixed-rank examples does not
establish that it works in general. Legality checks and the selected numerical
policy impose further requirements on the completed rewrite.

Our research question is: \emph{What if we could design a language that makes
tensor rewriting as easy and intuitive as writing its algebraic form, while
leaving complex attribute transformation to the compiler?}
TEPL pursues this goal by using operation semantics and symbolic axis mapping
to infer replacement attributes under explicit conditions.
Figure~\ref{fig:why-tepl} illustrates the idea with dot reassociation. It
compares a symbolic rewrite pattern, handwritten attribute logic, and two
ways to write the rule in TEPL: computing attributes explicitly or letting
the compiler infer them. The example has no batch axes and assumes an
arithmetic policy that permits reassociation. The completed rewrite must
still pass the applicability and semantic checks.

\paperfigure{1}{figures/Why TEPL_ From Algebraic Rules to Practical Tensor Rewrites.pdf}{TEPL combines concise tensor rewrite patterns with explicit or inferred replacement attributes.}{fig:why-tepl}{Four panels compare an algebraic pattern, custom implementation logic,
explicit TEPL attribute expressions, and inferred TEPL attributes for the
same nested-dot replacement graph.}
\FloatBarrier
\subsection{Central Thesis}
TEPL describes algebraic patterns, applicability conditions, and attribute
changes together. Users define operations once in reusable dialects, then
write replacement descriptors or request inference with \code{auto}.
The central technical advance pursued here is construction of reusable RHS
attribute programs from semantics and ordered coordinate correspondence.
Reasoning uses operation legality and declared conditions; inference does
not silently strengthen the guard when correctness cannot be established.

\subsection{Proposed Contributions}
\begin{enumerate}
\item \textbf{A language for tensor rewrites.} User-defined dialects express
operations and semantics; rules combine patterns, conditions, descriptor
capture, reuse, and explicit attribute computation.
\item \textbf{Automatic attribute inference.} For supported rewrites, the
compiler constructs reusable functions of source metadata without fixing
rank. Axes IR combines arbitrary-length groups with ordered interfaces,
operand mappings, and value semantics.
\item \textbf{Symbolic axis mapping for composed operations.} Compositional
mappings and group refinement preserve logical coordinate identity, order,
and pairings for read matching and attribute construction.
\end{enumerate}
These are design and mechanism claims, with the correctness contract stated
as an intended guarantee. Section~\ref{sec:evaluation} reports XLA-derived
coverage and conditional inference results; authoring and performance benefits
remain evaluation questions.
